\section{Current Situation}

\subsection{Replacing NASA with Private Companies}

The ongoing privatisation of space activities has impacted the public sector, such as NASA and international bodies such as the European Space Agency (ESA), which jointly coordinates the International Space Station (ISS) alongside several other national space agencies.\footnote{National Aeronautics and Space Administration. (n.d.). \emph{FAQs: The International Space Station transition plan.} NASA. \url{https://www.nasa.gov/faqs-the-international-space-station-transition-plan/}} The aging of the station pushes NASA and the White House to start working on its future successor, but aims to be a cost effective commercial approach to the low-orbit activities related to the ISS.\footnote{U.S. Government Accountability Office. (2025). \emph{Low-Earth orbit: NASA faces impending decisions for commercial transition} (GAO-26-107805). \url{https://www.gao.gov/products/gao-26-107805}} This allows NASA budget cuts with private backing of the project, allowing greater expansion in the future of low-orbit activities.\footnote{National Aeronautics and Space Administration. (n.d.). \emph{Commercial space stations}. NASA. \url{https://www.nasa.gov/humans-in-space/commercial-space/commercial-space-stations/}}

NASA has recently shifted its focus to the near future moon landing, projected to happen in 2028.\footnote{National Aeronautics and Space Administration. (n.d.). \emph{Moon to Mars: NASA's Artemis program}. NASA. \url{https://www.nasa.gov/humans-in-space/artemis/}} The SLS/Orion launcher already has ongoing live tests such as the recent Artemis 2 mission which brought astronauts around the moon to determine landing sites.\footnote{National Aeronautics and Space Administration. (2026, April 1). \emph{NASA's Artemis II mission leaves Earth orbit for flight around Moon} [Press release]. NASA. \url{https://www.nasa.gov/news-release/nasas-artemis-ii-mission-leaves-earth-orbit-for-flight-around-moon/}} This mission brought up debates on whether interplanetary exploration could be given to private entities such as Space X, the leader of the private space sector and its Falcon 9 launcher which is able to land back on Earth for reuse.\footnote{Roulette, J., \& Erwin, S. (2026, April 1). \emph{NASA's moon mission tests aerospace old guard as SpaceX, Blue Origin hover}. Reuters. \url{https://www.reuters.com/science/nasas-moon-mission-tests-aerospace-old-guard-spacex-blue-origin-hover-2026-04-01/}} Congress has so far refused proposals to terminate the Artemis program for commercial alternatives replacement and firmly rejected any attempts.\footnote{Berger, E. (2025, September). \emph{Congress and Trump may compromise on the SLS rocket by axing its costly upper stage}. Ars Technica. \url{https://arstechnica.com/space/2025/09/congress-and-trump-may-compromise-on-the-sls-rocket-by-axing-its-costly-upper-stage/}}

The House does support proposals of future permanent American outposts on the moon as soon as 2030, supported by crewed and cargo missions on regular schedules.\footnote{\emph{House appropriators advance key NASA funding bill}. (2026, May 13) The Planetary Society. \url{https://www.planetary.org/articles/house-appropriators-advance-key-nasa-funding-bill}} However, this plan would lean on commercial partnerships and supplemental resources backing.\footnote{Foust, J. (2026). \emph{House committee advances NASA authorization bill}. SpaceNews. \url{https://spacenews.com/house-committee-advances-nasa-authorization-bill/}}

\subsection{NASA Budget Struggles between Congress and the White House}

Since the arrival of Donald Trump in the White House in late 2024, NASA has faced multiple attempts of budget cutting by the 47th presidential administration, which has led to multiple proposals to Congress to allow more power to private companies on space activities.\footnote{Foust, J. (2025). \emph{White House budget proposal would phase out SLS and Orion, scale back ISS operations}. SpaceNews. \url{https://spacenews.com/white-house-budget-proposal-would-phase-out-sls-and-orion-scale-back-iss-operations/}} The latest proposal comes from the FY2026 White House proposal to cut 24\% of NASA’s budget from \$24.8 billion to \$18.8 billion, the lowest since 2015, also including a 47\% cut of its science program, one of the most important ones in terms of space research and progress.\footnote{Wall, M. (2025, May). \emph{Trump's 2026 budget would slash NASA funding by 24\% and its workforce by nearly one-third}. Space.com. \url{https://www.space.com/space-exploration/trumps-2026-budget-would-slash-nasa-funding-by-24-percent-and-its-workforce-by-nearly-one-third}} However, the Senate and Congress jointly rejected the proposal, instead settling with the administration to a \$400 million cut from NASA’s prior budget.\footnote{American Astronomical Society. (2026, January). \emph{Congress passes fiscal year 2026 spending bills for NSF, NASA, and DOE}. AAS. \url{https://aas.org/posts/news/2026/01/congress-passes-fiscal-year-2026-spending-bills-nsf-nasa-and-doe}}

For FY2027, the same scenario is repeated, with Trump proposing a 23\% spending cut for 2027.\footnote{Foust, J. (2026, April). \emph{White House again proposes steep NASA budget cuts}. SpaceNews. \url{https://spacenews.com/white-house-again-proposes-steep-nasa-budget-cuts/}} This proposal is already faced with a lot of criticism, such as Rep. Brian Babin, chairman of the House Committee on Science, Space and Technology called this proposal “dead on arrival”.\footnote{\emph{House Science Committee members vow to reject NASA budget cuts}. (2026). Aerospace America (AIAA). \url{https://aerospaceamerica.aiaa.org/house-science-committee-members-vow-to-reject-nasa-budget-cuts/}} The House has advanced a bill cutting flat the budget to \$24.4 billion, rejecting the White House’s proposal while cutting the Science Mission Directorate by 17\%.\footnote{\emph{House appropriators advance key NASA funding bill}. (2026, May 13). The Planetary Society. \url{https://www.planetary.org/articles/house-appropriators-advance-key-nasa-funding-bill}}

\subsection{Planned bills and loopholes}

While the majority of the area regarding space activities has been regulated by laws, some small loopholes remain unaddressed by the U.S. government to this day. For example, activities such as on-orbit servicing, in-space manufacturing or even debris removal do not fall under any regulated activities or any agency supervision.\footnote{Belfer Center for Science and International Affairs. (2025). \emph{Explainer: Novel space activities in the executive order on enabling competition in the commercial space industry}. Harvard Kennedy School. \url{https://www.belfercenter.org/research-analysis/explainer-novel-space-activities-executive-order-enabling-competition-commercial}} Although some laws have been or will be passed directly by the White House. Executive Order 14335 from August 2025 enables competition in the commercial space industry, enabling future reforms in major areas such as commercial launch and reentry, spaceport infrastructure, novel space activity authorisation and leadership and accountability, aiming for an increasing commercial launch schedule by 2030.\footnote{The White House. (2025, August 13). \emph{Enabling competition in the commercial space industry} [Presidential action]. \url{https://www.whitehouse.gov/presidential-actions/2025/08/enabling-competition-in-the-commercial-space-industry/}} The administration also signed an executive order in December 2025, “Ensuring American Space Superiority,” aimed at future lunar activities such as returning astronauts by 2028 as well as deploying a lunar outpost and a nuclear reactor on its surface by 2030.\footnote{The White House. (2025, December 18). \emph{Ensuring American space superiority} [Presidential action]. \url{https://www.whitehouse.gov/presidential-actions/2025/12/ensuring-american-space-superiority/}}

\subsection{Key takeaways}

The House has mostly acted as a counterweight to the white house’s private proposals regarding NASA budget and takeover, overly considered as aggressive proposals. Nevertheless, congress is still considering and building a legal and regulatory framework regarding commercial space activities and private space stations. Branches of the executive government still push for companies such as Space X to be favored based on experience, while congress desires to maintain public space and scientific programs, ensuring NASA’s sovereignty in the future.
